% ===== §23 The Beauty of Transmission =====
\section{The Beauty of Transmission: Why Aesthetic Education Is an Education}
\label{p22:sec:transmission}

\noindent
The paper has described what aesthetic education cultivates and how, but it has deferred the question lodged in its own name, the question of the word education, of the cultivating of one by another, the passing of the capacity from those who have it to those who do not. The four-step model described how a person who already has the sensibility exercises it; the worked example described how that person trains the capacity in herself. Neither has yet said how the sensibility passes from one person to another, how the capacity to build generative fields is transmitted, which is what the word education finally demands. This closing section addresses the transmission, and it finds that the manner of the transmission is dictated by everything the paper has established, so that aesthetic education must be an education in a particular and demanding sense, a transmission that cannot be reduced to instruction, and that the paper itself, in being read, is one instance of the transmission it describes.

\subsection{The limit of instruction in transmission}
\label{p22:sub:trans-not-instruction}

The transmission of the sensibility cannot be the transmission of a code, for the reason the whole paper has pressed: the sensibility is a trained perception of the state-dependent appropriate, and the state-dependent appropriate cannot be specified, so there is no code to transmit. One cannot teach the capacity to build generative fields by telling, by handing over rules for the appropriate gesture, because the appropriate gesture is a property of the token field and no rule reaches tokens. This is why aesthetic education is an education and not an instruction, and the distinction, drawn several times in the paper, can now be given its full weight. An instruction transmits a specifiable content, a code that the instructed can apply; an education raises a capacity that the educated must exercise in the concrete, and it raises it not by transmitting a content but by cultivating a power. The transmission of the sensibility is education in this strict sense: it cannot hand over what it transmits, because what it transmits is a perception, and a perception is not the kind of thing that can be handed over, only the kind of thing that can be cultivated in another by exercise.

\subsection{Transmission requires co-presence}
\label{p22:sub:trans-copresence}

If the sensibility cannot be transmitted by instruction, how is it transmitted at all? The answer the paper's framework dictates is that it is transmitted by co-presence, by the shared exercise of the capacity in which one who has it and one who is acquiring it perceive, receive, share, and reproduce together, so that the acquiring one's perception is cultivated by participation in the having one's. The sensibility passes as the capacity for joint attention passes, by being exercised jointly, the finer perception drawing the coarser one along, the way a child's perception of a field is cultivated by sharing fields with those who perceive them finely, the way any trained perception is acquired by exercising it alongside one already trained. This is transmission by demonstration and shared practice rather than by telling, and it requires the co-presence of the transmitter and the acquirer in a shared field, because only in a shared field can the acquiring perception be drawn along by the having one. The sensibility cannot be sent; it can only be grown, in a shared field, by joint exercise, and this is the deepest reason aesthetic education is an education, that its transmission has the structure of the very sharing the paper has described, the joint attention in which a field is built between two.

\claim{The sensibility is transmitted neither by instruction nor by any sending of content, but by co-presence, the shared exercise of the capacity in which an acquiring perception is cultivated by joint participation in a finer one. Transmission has the structure of sharing itself, the joint attention in which a field is built between two, so that to transmit the capacity to build generative fields is to build a generative field with the one who is acquiring it. Aesthetic education is an education in the strictest sense, a cultivation of perception by shared exercise, and its transmission cannot be reduced to instruction because a perception cannot be handed over, only grown in a shared field.}

This gives the word education its full sense and closes the question the section opened. The cultivating that the word names is the cultivation of one's perception by another's, in a shared field, by joint exercise, and it is therefore not a lesser or metaphorical education but education in its most exact form, the growing of a capacity in another by sharing the exercise of it. That the transmission requires co-presence, a shared field in which the capacities are jointly exercised, is not a limitation but the mark of the kind of thing being transmitted, a perception that lives only in its exercise and can be cultivated only by being exercised together.

\subsection{The child and the inheritance of a sensibility}
\label{p22:sub:trans-child}

The clearest case of the transmission is the one in which a child acquires the sensibility from those who raise her, and it is worth dwelling on because it shows the structure of the transmission at its most consequential and confirms that the transmission cannot be instruction. A child does not learn to build generative fields by being told the rules for building them, because there are no such rules, and a child raised on instruction alone, told what to do and not shown, acquires no sensibility at all. What cultivates the child's perception is participation in fields built by those who can build them, the daily sharing of a life in which the contingencies are caught and shared and the daily things renewed, so that the child's own perception is drawn along by the finer perception she lives within, acquiring by exercise what no telling could convey. A child raised in a home whose fields are generative acquires, without a lesson being given, the perception of what makes a field generative, the way a child raised among a language's speakers acquires the language without being taught its grammar, by exercising the capacity alongside those who have it. And a child raised in a home whose fields are flattened or captured acquires the coarser perception, or the defended one, not by being taught it but by living within it, which is the somber side of the same truth, that the sensibility is transmitted by the fields one is raised within and not by the precepts one is given.

This has a consequence that the political economy section's concern with the dispossessed must be allowed to darken. If the sensibility is transmitted by participation in generative fields, then those raised in homes where the conditions of field-building have been damaged, by exhaustion, by captured yield, by the manufactured pressures that flatten the everyday, inherit a diminished capacity through no failure of their own, the structural injustice reproducing itself in the transmission of the sensibility across generations. The cultivation this paper describes is available to the dispossessed, as the political economy argued, requiring no surplus they lack; but the transmission of the capacity to the next generation is harder where the fields the child is raised within have been damaged by the conditions the dispossessed endure, and this is a further face of the structural remainder, the part the cultivation alone cannot reach. The honest account holds both: the sensibility is transmissible in any home by the sharing of fields, requiring no surplus; and the damage that captured and exhausting conditions do to a home's fields is reproduced in what its children inherit, so that the justice of the field bears not only on those who build it now but on the capacity their children will have to build it in their turn.

\subsection{The paper, an instance of its own transmission}
\label{p22:sub:trans-self-instance}

The paper closes by turning upon itself, as the series has turned its papers upon themselves throughout, and finding that it is an instance of the transmission it describes. A paper cannot, by the argument just made, transmit the sensibility by instruction, and this paper has not tried to, has offered no code for the appropriate gesture, no rules for building fields, because it has held throughout that no such code exists. What it has tried to do is otherwise, and the manner of its trying is the manner of the transmission it describes. It has cultivated, in the reader who has followed it, a perception, drawing the reader's attention to the state-dependence of the appropriate, to the difference between the spiral and the wheel, to the failures that shadow each movement of the cycle, in the hope that the reader's own perception of the fields of her own life is thereby a little finer, a little more able to discern what before went unseen. This cultivation is not instruction but a kind of shared exercise conducted in the medium of reading, the paper perceiving the structure of the field and the reader drawn along to perceive it too, and to the extent it has succeeded the reader's sensibility has been cultivated by the paper's, which is the transmission the paper describes performed in the only co-presence a paper allows.

This self-instancing is not a flourish but a consequence, and it completes the paper's form in the way the series has always sought, the form enacting the claim. The paper claims that the sensibility is transmitted by the cultivation of one perception by another in a shared exercise; and the paper, in being read, cultivates the reader's perception by its own, in the shared exercise of attending together to the structure of the field. It cannot do this by instruction, and it has not; it has done it, to whatever extent it has done it, by drawing the reader along into a finer perceiving, which is exactly the transmission it describes. The paper is thus its own last example, a field built between writer and reader in the medium of the page, accumulating, if it has worked, a perception in the reader that was not there before, and returning, at its close, to the root from which it began, the cultivation of the capacity to build generative fields, now performed upon the reader rather than only described to her. If the reader perceives the fields of her life a little more finely for having read, the paper has built a field with her, and its argument has become, in the reading, the thing it argued for.
